% Generated draft: requires mathematical and source review.
\documentclass[11pt]{article}
\usepackage[margin=25mm]{geometry}
\usepackage{fontspec}
\setmainfont{Latin Modern Roman}
\usepackage{amsmath,amssymb,mathtools,unicode-math}
\setmathfont{Latin Modern Math}
\usepackage{enumitem,needspace,xurl,hyperref}
\setlength{\parindent}{0pt}
\setlength{\parskip}{5pt}
\setlength{\emergencystretch}{3em}
\newcommand{\N}{\mathbb{N}}
\newcommand{\Z}{\mathbb{Z}}
\newcommand{\Q}{\mathbb{Q}}
\newcommand{\R}{\mathbb{R}}
\newcommand{\C}{\mathbb{C}}
\newcommand{\F}{\mathbb{F}}
\newcommand{\A}{\mathbb{A}}
\newcommand{\PP}{\mathbb{P}}
\newcommand{\E}{\mathbb{E}}
\newcommand{\eps}{\varepsilon}
\newcommand{\dd}{\,\mathrm{d}}
\begin{document}
\section*{UnsolvedMath ID: 6200050}
\textbf{Source problem code:} catalogue entry
\subsection{Definitions and mathematical statement}
The initial request, whether a prior construction can use smaller \(m\), specifically \(m=n+1\), has no antecedent in the extracted record and cannot be formalized independently. The text also says that the same problem applies to complex hyperbolic space.

The explicit Kapovich subproblem concerns a standard inclusion of hyperbolic symmetric spaces \(W_n\hookrightarrow W_{n+1}\) and its boundary map
\[
i:X=\partial_\infty W_n\hookrightarrow Y=\partial_\infty W_{n+1}.
\]
The extracted symbol ``HHn'' is not defined. If it denotes quaternionic hyperbolic space \(\mathbb H_{\mathbb H}^{n}\), that choice and its admissible dimension range must be confirmed; a real- or complex-hyperbolic interpretation is a different problem. The boundaries are equipped with the intended visual metric structures.

A quasisymmetric embedding \(f:X\to Y\) is an embedding onto its image with a distortion homeomorphism \(\eta:[0,\infty)\to[0,\infty)\) such that
\[
\frac{d_Y(f(x),f(a))}{d_Y(f(x),f(b))}
\le\eta\left(\frac{d_X(x,a)}{d_X(x,b)}\right)
\]
for all distinct \(x,a,b\in X\). Local rigidity asks whether every such \(f\) sufficiently close to \(i\) has the form
\[
f=(\partial_\infty g)\circ i
\qquad\text{for some }g\in\operatorname{Isom}(W_{n+1}).
\]
The equality concerns the parametrized embedding, not merely congruence of its image. Since \(X\ne Y\), the source's identity means the inclusion \(i\).

For example, in the compact-open boundary topology a precise version would require some \(\varepsilon>0\) such that
\[
\sup_{x\in X}d_Y(f(x),i(x))<\varepsilon
\quad\Longrightarrow\quad f=(\partial_\infty g)\circ i.
\]
The source does not specify this topology, metric normalization, or whether a fixed quasisymmetry control \(\eta\) is required. The displayed neighborhood is therefore a candidate convention, not an added verified hypothesis. The supplied research summary calls the specific boundary rigidity statement open and mentions related higher-rank embedding rigidity without establishing this conclusion.
\subsection{Short English statement}
Is the standard boundary inclusion locally rigid among quasisymmetric embeddings, so every sufficiently nearby embedding is induced by an ambient isometry? The hyperbolic-space notation, closeness topology, and missing smaller-dimension antecedent need source confirmation.
\subsection{Sources}
\begin{itemize}
\item Local UnsolvedMath registry, permanent ID \(6200050\), exact source code \texttt{catalogue entry}, ``Boundaries of Groups and Kleinian Groups --- Problem 50.'' Proposer David Fisher, 2005; explicit subproblem attributed to Misha Kapovich.
\item Kapovich, \emph{Problems on Boundaries of Groups and Kleinian Groups} (2005), Problem 50, recorded PDF page 15; \url{https://www.math.ucdavis.edu/~kapovich/EPR/problems.pdf}. Recorded access 2026-07-29; pdf-text-normalized extraction.
\item Supplied OPEN-TRIAGE research summary mentions Kleiner--Leeb and Eskin--Fisher--Whyte as related rigidity literature, without complete references or a theorem resolving this boundary statement.
\end{itemize}
\end{document}
